
\section*{练习}
\label{sec:ch03-exercises}

\vkexsync{13}

\begin{exercise}\label{ex:ch03-13}
与\cref{ex:ch02-12}中所定义的函数以及同一练习中所建议的那个\trVectorField{}相伴的\trComplex，其\trHomology{}是什么？
\end{exercise}

\begin{exercise}\label{ex:ch03-14}
设 $E$ 与 $F$ 是一个有限维实\trVectorSpace{}的两个向量子空间。证明：$E$ 的一个\trOrientation，就是 $E$ 的诸基在等价关系
\[
\mathcal{B}\sim\mathcal{B}'\Longleftrightarrow \det_{\mathcal{B}}\mathcal{B}'>0
\]
之下的一个等价类。

同样地，验证：关系
\[
\mathcal{B}\sim\mathcal{B}'\Longleftrightarrow \det_{(\mathcal{B},\mathcal{B}_0)}(\mathcal{B}',\mathcal{B}_0)>0
\]
在 $F$ 的诸\trComplement{}的诸基上定义了一个等价关系，它不依赖于所选定的 $F$ 的基 $\mathcal{B}_0$。这些等价类就是 $F$ 的诸\trCoorientation。

再验证：若 $E$ 已\trOrientation、$F$ 已\trCoorientation，并且 $E$ 与 $F$ \trTransverse，则 $E\cap F$ 便被赋以一个\trOrientation。\footnote{译注：原书末词印作 co-oriented（\trCoorientation）；按 $E$ 与 $F$ \trTransverse、$E+F=W$，\trExact{}列 $0\to E\cap F\to E\to W/F\to 0$ 连同 $E$ 的\trOrientation{}与 $W/F$（借 $F$）的\trOrientation，给出的乃是 $E\cap F$ 的\trOrientation{}而非其在 $W$ 中的\trCoorientation（反例：$W=\RR^3$、$E=\operatorname{span}(e_1,e_2)$、$F=\operatorname{span}(e_2,e_3)$ 以 $e_1$ 为其\trCoorientation，则 $\operatorname{diag}(1,1,-1)$ 保持全部已知数据却反转 $W/(E\cap F)$ 的\trOrientation），疑为原书笔误；此处径改。——译者注}
\end{exercise}

\begin{exercise}\label{ex:ch03-15}
对于本书中所用的诸\trManifold{} $\CP^n(\CC)$、$T^2$ 与 $S^n$ 上的那些\trMorseFunction{}例子，确定\trComplex{} $\bigl(C_{\star}(f;\ZZ),\partial_X\bigr)$ 的\trHomology。
\end{exercise}
